% ═══════════════════════════════════════════════════════════════════════════
\section{Listati Python completi}\label{app:listati}
% ═══════════════════════════════════════════════════════════════════════════

In questa appendice sono riportati i listati Python completi dei tre algoritmi
discussi nel Capitolo~\ref{sec:algoritmi} e implementati nell'applicazione
interattiva (Capitolo~\ref{sec:esperimenti}). I listati fanno riferimento alle
seguenti funzioni ausiliarie, definite dall'utente e dipendenti dal dataset:

\begin{itemize}
  \item \texttt{loss\_i(w, i)}: perdita individuale $\ell(w;i)$,
  \item \texttt{grad\_i(w, i)}: gradiente individuale $\nabla\ell(w;i)$,
  \item \texttt{hessvec\_i(w, i, v)}: prodotto $\nabla^2\ell(w;i)\,v$,
  \item \texttt{grad\_full(w)}: gradiente esatto $\nabla J(w)$ (usato solo per
  il criterio di arresto),
  \item $N$: dimensione del dataset.
\end{itemize}

\subsection{Listato B.1: metodo del gradiente a campione dinamico}

\begin{pythoncode}[title= Dynamic GD]

import numpy as np

def dynamic_gd(w0, theta, max_iter, alpha, batch0):
    w = np.array(w0, dtype=float)
    n = max(batch0, 2)
    history     = [w.copy().tolist()]
    batch_sizes = [n]

    for k in range(max_iter):
        indices = np.random.choice(N, size=n, replace=False)
        grads = np.array([grad_i(w, i) for i in indices])
        g = np.mean(grads, axis=0)

        # Condizione di Controllo della Varianza (CCV)
        if n > 1:
            var_vec = np.var(grads, axis=0, ddof=1)
        else:
            var_vec = np.zeros_like(g)
        V_norm1 = np.sum(var_vec)

        gg = np.dot(g, g)
        if gg > 1e-16:
            if V_norm1 / n > theta**2 * gg:
                n_new = int(np.ceil(V_norm1 / (theta**2 * gg))) + 1
                n = min(n_new, N)

        # Line search di Wolfe sul batch corrente
        def J_batch(w_curr):
            return np.mean([loss_i(w_curr, i) for i in indices])

        c1, c2 = 1e-4, 0.9
        step = alpha
        J_curr = J_batch(w)
        g_norm2 = np.dot(g, g)
        d = -g
        gd = -g_norm2
        if g_norm2 > 1e-16:
            for _ in range(30):
                w_new = w + step * d
                if J_batch(w_new) <= J_curr + c1 * step * gd:
                    g_new = np.mean([grad_i(w_new, i) for i in indices], axis=0)
                    if np.dot(g_new, d) >= c2 * gd:
                        break
                step *= 0.5
            else:
                step = 0.0

        w = w + step * d

        if np.linalg.norm(grad_full(w)) < 1e-6:
            history.append(w.copy().tolist())
            batch_sizes.append(n)
            break

        history.append(w.copy().tolist())
        batch_sizes.append(n)

    return history, batch_sizes

history, batch_sizes = dynamic_gd(w0, theta, max_iter, alpha, batch0)
\end{pythoncode}


\subsection{Listato B.2: Newton-CG con campionamento dinamico}

\begin{pythoncode}[title= Newton-CG]

import numpy as np

def cg(A, b, gamma, maxcg):
    """Gradiente coniugato per risolvere A x = b (A operatore Hessiana-vettore)."""
    x = np.zeros_like(b)
    r = b - A(x)
    p = r.copy()
    rr = np.dot(r, r)
    for _ in range(maxcg):
        Ap = A(p)
        pHp = np.dot(p, Ap)
        if pHp <= 1e-14:
            break
        alpha = rr / pHp
        x = x + alpha * p
        r_new = r - alpha * Ap
        rr_new = np.dot(r_new, r_new)
        # Test di arresto adattivo: residuo confrontabile con la soglia
        if rr_new <= gamma * np.dot(x, x) + 1e-16:
            return x
        beta = rr_new / rr
        p = r_new + beta * p
        r = r_new
        rr = rr_new
    return x

def newton_cg(w0, theta, max_iter, alpha, batch0, R, maxcg):
    w = np.array(w0, dtype=float)
    n = max(batch0, 2)
    history, batch_sizes = [w.copy().tolist()], [n]

    for k in range(max_iter):
        indices_S = np.random.choice(N, size=n, replace=False)
        g = np.mean([grad_i(w, i) for i in indices_S], axis=0)

        # Sottocampione per l'Hessiana
        n_h = min(max(1, int(round(R * n))), N)
        indices_H = np.random.choice(indices_S, size=n_h, replace=False)
        Hv = lambda v: np.mean([hessvec_i(w, i, v) for i in indices_H], axis=0)

        # Coefficiente gamma: varianza dei prodotti Hessiana-vettore lungo p0
        p0 = -g
        p0_norm2 = np.dot(p0, p0)
        gamma = 0.0
        if p0_norm2 > 1e-16 and n_h > 1:
            Hp0 = np.array([hessvec_i(w, i, p0) for i in indices_H])
            gamma = np.sum(np.var(Hp0, axis=0, ddof=1)) / (n_h * p0_norm2)

        d = cg(Hv, -g, gamma, maxcg)

        # Line search di Wolfe lungo d
        J_batch = lambda wc: np.mean([loss_i(wc, i) for i in indices_S])
        c1, c2 = 1e-4, 0.9
        step, J_w = alpha, J_batch(w)
        gd = np.dot(g, d)
        if gd >= 0:
            d = -g
            gd = -np.dot(g, g)
        for _ in range(30):
            w_new = w + step * d
            if J_batch(w_new) <= J_w + c1 * step * gd:
                g_new = np.mean([grad_i(w_new, i) for i in indices_S], axis=0)
                if np.dot(g_new, d) >= c2 * gd:
                    break
            step *= 0.5
        else:
            step = 0.0
        w = w + step * d

        history.append(w.copy().tolist())
        batch_sizes.append(n)

        # CCV sul nuovo campione
        indices_new = np.random.choice(N, size=n, replace=False)
        g_new = np.mean([grad_i(w, i) for i in indices_new], axis=0)
        var_vec = (np.var([grad_i(w, i) for i in indices_new], axis=0, ddof=1)
                   if n > 1 else np.zeros_like(g_new))
        V_norm1, gg_new = np.sum(var_vec), np.dot(g_new, g_new)
        if gg_new > 1e-16 and V_norm1 / n > theta**2 * gg_new:
            n = min(int(np.ceil(V_norm1 / (theta**2 * gg_new))) + 1, N)

        if np.linalg.norm(grad_full(w)) < 1e-6:
            break
    return history, batch_sizes

history, batch_sizes = newton_cg(w0, theta, max_iter, alpha, batch0, R, maxcg)
\end{pythoncode}



\subsection{Listato B.3: Newton-CG per problemi $L_1$-regolarizzati}

\begin{pythoncode}[title= Newton-L1]

import numpy as np

def newton_l1(w0, theta, max_iter, alpha, batch0, nu, sigma, maxcg,
              eta=0.5, R=0.1):
    w = np.array(w0, dtype=float)
    n = max(batch0, 2)
    history     = [w.copy().tolist()]
    batch_sizes = [n]

    # F_batch: obiettivo regolarizzato sul batch
    def F_batch(v, indices):
        Jb = np.mean([loss_i(v, i) for i in indices])
        return Jb + nu * np.sum(np.abs(v))

    # subgrad_batch: gradiente generalizzato (6.3)
    def subgrad_batch(v, indices):
        grads = np.array([grad_i(v, i) for i in indices])
        gJ = np.mean(grads, axis=0)
        g = np.zeros_like(v)
        for i in range(len(v)):
            if v[i] > 0:
                g[i] = gJ[i] + nu
            elif v[i] < 0:
                g[i] = gJ[i] - nu
            else:
                if gJ[i] < -nu:
                    g[i] = gJ[i] + nu
                elif gJ[i] > nu:
                    g[i] = gJ[i] - nu
                else:
                    g[i] = 0.0
        return g

    # project_orthant: proiezione sulla faccia ortante Omega_k (6.5)
    def project_orthant(v, z):
        res = v.copy()
        for i in range(len(v)):
            if z[i] != 0 and np.sign(res[i]) != z[i]:
                res[i] = 0.0
        return res

    for k in range(max_iter):
        indices_S = np.random.choice(N, size=n, replace=False)
        grads = np.array([grad_i(w, i) for i in indices_S])
        g_batch = np.mean(grads, axis=0)

        # vettore z_k che caratterizza la faccia ortante (6.4)
        z = np.where(w > 0, 1,
            np.where(w < 0, -1,
                np.where(g_batch < -nu, 1,
                    np.where(g_batch > nu, -1, 0))))

        sg = subgrad_batch(w, indices_S)
        sgn = np.linalg.norm(sg)
        if sgn < 1e-10:
            history.append(w.copy().tolist())
            batch_sizes.append(n)
            break

        # Hessiana sottocampionata sulle variabili libere
        n_h = min(max(1, int(round(R * n))), N)
        indices_H = np.random.choice(indices_S, size=n_h, replace=False)
        free = (z != 0)
        d = np.zeros_like(w)
        if np.any(free):
            g_free = sg[free]

            def Hv(v_full):
                return np.mean([hessvec_i(w, i, v_full) for i in indices_H],
                               axis=0)

            def Hv_free(v_free):
                v_full = np.zeros_like(w)
                v_full[free] = v_free
                return Hv(v_full)[free]

            # CG Hessian-free sul sottospazio libero
            tol_cg = eta * np.linalg.norm(g_free)
            d_free = np.zeros(np.sum(free))
            r = -g_free.copy()
            p = r.copy()
            rr = np.dot(r, r)
            for _ in range(maxcg):
                Hp = Hv_free(p)
                pHp = np.dot(p, Hp)
                if pHp <= 1e-14:
                    if np.linalg.norm(d_free) < 1e-14:
                        d_free = -g_free.copy()
                    break
                alpha_cg = rr / pHp
                d_free = d_free + alpha_cg * p
                r_new = r - alpha_cg * Hp
                rr_new = np.dot(r_new, r_new)
                if np.sqrt(rr_new) <= tol_cg:
                    break
                beta = rr_new / rr
                p = r_new + beta * p
                r = r_new
                rr = rr_new
            d[free] = d_free

        # Ricerca lineare proiettata (Armijo con backtracking)
        step = alpha
        F_w = F_batch(w, indices_S)
        sg_d = np.dot(sg, d)
        if sg_d >= 0:
            d = -sg
            sg_d = -np.dot(sg, sg)
        w_new = w.copy()
        for _ in range(20):
            w_trial = project_orthant(w + step * d, z)
            if F_batch(w_trial, indices_S) <= F_w + sigma * step * sg_d:
                w_new = w_trial
                break
            step *= 0.5
            if step < 1e-12:
                w_new = w.copy()
                break
        w = w_new

        # CCV estesa al gradiente della parte liscia
        indices_new = np.random.choice(N, size=n, replace=False)
        grads_new = np.array([grad_i(w, i) for i in indices_new])
        g_new = np.mean(grads_new, axis=0)
        if n > 1:
            var_vec = np.var(grads_new, axis=0, ddof=1)
        else:
            var_vec = np.zeros_like(g_new)
        V_norm1 = np.sum(var_vec)
        gg_new = np.dot(g_new, g_new)
        if gg_new > 1e-16 and V_norm1 / n > theta**2 * gg_new:
            n_new = int(np.ceil(V_norm1 / (theta**2 * gg_new))) + 1
            n = min(n_new, N)

        history.append(w.copy().tolist())
        batch_sizes.append(n)
        if np.linalg.norm(grad_full(w)) < 1e-6:
            break
    return history, batch_sizes

history, batch_sizes = newton_l1(w0, theta, max_iter, alpha, batch0,
                                 nu, sigma, maxcg, eta)
\end{pythoncode}

